\documentclass[a4paper,12pt]{article}
\usepackage{amsmath}
\usepackage{geometry}
\geometry{margin=1in}
\usepackage{booktabs}
\usepackage{parskip}
\usepackage[T1]{fontenc}
\usepackage{times}

\begin{document}

\title{TrajectoryModel Architecture Description}
\author{}
\date{}
\maketitle

\section{Overview}
The TrajectoryModel predicts future vehicle trajectories from two RGB images (current and previous frames). It uses a motion-aware encoder that extracts multi-scale features from each frame, fuses them with temporal differences, and then applies a lightweight transformer decoder to predict the future displacement vectors. An optional auxiliary head can predict a 2D dynamics vector from pooled features. The default input resolution is (B, 3, 360, 640), and the default output is (B, 12, 2) with non-uniform time offsets spanning 3 seconds.

\section{Model Components}

\subsection{Utility Blocks}
\begin{itemize}
    \item \textbf{ConvGNAct}: A 2D convolution followed by group normalization and SiLU activation.
    \item \textbf{SEBlock}: Squeeze-and-Excitation block that rescales channels using a global pooling and two-layer MLP.
    \item \textbf{Residual}: Two ConvGNAct layers with a residual connection, optionally followed by SE attention.
\end{itemize}

\subsection{MotionBackbone}
A shared-weight CNN extracts features from each frame at two spatial scales.
\begin{itemize}
    \item \textbf{Input}: Image of shape (B, 3, H, W).
    \item \textbf{Architecture}:
        \begin{itemize}
            \item \textit{Stem}: ConvGNAct (3 to base, kernel=5, stride=2), Residual (base, SE enabled).
            \item \textit{Stage1}: ConvGNAct (base to 2\,base, stride=2), Residual (2\,base, SE enabled).
            \item \textit{Stage2}: ConvGNAct (2\,base to 3\,base, stride=2), Residual (3\,base, SE enabled). Output is 1/8 resolution.
            \item \textit{Stage3}: ConvGNAct (3\,base to 4\,base, stride=2), Residual (4\,base, SE enabled). Output is 1/16 resolution.
        \end{itemize}
    \item \textbf{Output}: Feature maps at 1/8 and 1/16 resolution.
\end{itemize}

\subsection{MotionFpnEncoder}
Fuses motion cues by comparing features between the current and previous frames.
\begin{itemize}
    \item \textbf{Per-scale fusion}: Concatenate current features, previous features, and their difference, then apply a 1x1 convolution to a shared feature dimension.
    \item \textbf{Top-down fusion}: Upsample the 1/16 features and add to the 1/8 features, followed by a ConvGNAct fusion block.
    \item \textbf{Coordinate channels}: Append normalized x/y coordinate maps in the range [-1, 1].
    \item \textbf{Output}: Feature map of shape (B, feat\_dim + 2, H/8, W/8). For 360x640 input, this is 45x80.
\end{itemize}

\subsection{VectorTransformerDecoder}
The decoder predicts displacement vectors in parallel using a transformer decoder.
\begin{itemize}
    \item \textbf{Memory}: The encoder feature map is projected with a 1x1 convolution and flattened to a sequence of length H\,W, followed by LayerNorm.
    \item \textbf{Queries}: A learned query embedding of length pred\_steps is added to a time-step MLP embedding.
    \item \textbf{Transformer decoder}: A lightweight decoder (default 2 layers, 4 heads) attends to the spatial memory.
    \item \textbf{Output head}: Two-layer MLP that predicts (x, y) per step.
    \item \textbf{Output}: Predicted trajectories of shape (B, pred\_steps, 2). Attention maps are not explicitly returned (filled with None).
\end{itemize}

\subsection{TrajectoryModel}
The full model integrates the encoder, decoder, and optional auxiliary head.
\begin{itemize}
    \item \textbf{Parameters}: Feature dimension (feat\_dim), hidden dimension (hidden\_dim), prediction steps (pred\_steps=12), optional auxiliary dynamics (default False).
    \item \textbf{Components}:
        \begin{itemize}
            \item MotionFpnEncoder (outputs feat\_dim + 2 channels with coordinate maps).
            \item VectorTransformerDecoder (uses feat\_dim + 2).
            \item Optional auxiliary head: AdaptiveAvgPool2d -> MLP (feat\_dim + 2 to 128 to 2).
        \end{itemize}
    \item \textbf{Time offsets}: If not provided, non-uniform offsets are generated for a 3-second horizon (0.1s x 6, 0.25s x 4, 0.7s x 2).
    \item \textbf{Forward}: Returns predicted trajectories, optional auxiliary output, and a placeholder attention list.
\end{itemize}

\section{Summary}
The TrajectoryModel combines motion-aware multi-scale encoding with a transformer decoder to predict future vehicle trajectories from image pairs. The design is lightweight, supports non-uniform time steps, and optionally produces a compact auxiliary dynamics output.

\end{document}